\chapterpage{T02 AND T04\enspace /\enspace SEPTEMBER CHECKPOINTS}{Local output counts and real input growth}
\subhead{T02\enspace /\enspace The fixed quartic}
For $P_4(x,y)=(x+2y)^4+3x-y$, the September checkpoint counts outputs at every precision in a specific local ring. In the local cubic extension used there, $7\mathcal O=\mathfrak p^3$ and the residue field has seven elements. Take $R_m=\mathcal O/\mathfrak p^{3m}$ and the full critical residue class $E_m=(\alpha+\mathfrak p)\bmod\mathfrak p^{3m}$ with $\alpha^3\equiv-1\pmod{\mathfrak p}$. The C4 support count is
\[
 |P_4(E_m,E_m)|=
 \begin{cases}
  4,&m=1,\\
  22\cdot7^{3m-4},&m\ge2.
 \end{cases}
\]
For $m\ge2$ the output support consists of three full $\mathfrak p^3$-cosets and one full $\mathfrak p^4$-coset. The relative density improves once. After both input restrictions are charged, its factor stabilizes at $22/49$. The proposed real full-interval packing must also pay for the number of inputs and the widths of the output classes. The recorded certificate fails to produce a positive net saving. Other higher-precision subsets remain outside that conclusion. \src{E7}

\subhead{T04\enspace /\enspace The minus-sign cubic}
The fourth track asks about $P_-(x,y)=(x-y)^3-3(x+y)$. Its two real input maps have opposite shifts and reciprocal positive slopes. Let $\epsilon=2+\sqrt3$, $a_0=\sqrt3/2$ and set
\[
 A_0=\{0\},\qquad
 A_{m+1}=(\epsilon A_m+a_0)\cup(\epsilon^{-1}A_m-a_0).
\]
The $2^m$ histories need not produce different scalars. The input theorem counts injective subfamilies and controls the number of distinct affine maps. For $\Lambda>1$ satisfying $\Lambda^{-1}+\Lambda^{-3}+\Lambda^{-5}=1$, and $\varphi=(1+\sqrt5)/2$, it gives
\[
 |A_m|=\exp(a_*m+o(m)),\qquad \log\Lambda\le a_*\le\log\varphi .
\]
All scalars lie in a fixed quadratic field with bounded denominator. Mixed branch pairs still move both the sum and the difference of the endpoints. The full output support has only the certified upper bound $\exp(2a_*m+o(m))$. The true exponential gap is unknown. No T01 plus-sign result supplies this estimate. \src{E9}

\smallskip\tagword{Evidence status.} The two results above are reported from identified checkpoint proofs. Both are summarized here. The dated status notes distinguish recovered copies, corrections and still-missing earlier provenance.

\chapterpage{T03 AND SHARED WORK\enspace /\enspace SEPTEMBER--NOVEMBER}{Norm symmetry and the cost of a return}
\subhead{T03\enspace /\enspace The Eisenstein quadratic}
For $Q(x,y)=x^2+xy+y^2$, write $\omega^2+\omega+1=0$. The product $(x-\omega y)(x-\omega^2y)=Q(x,y)$ identifies the output with a norm. Multiplying $z=x-\omega y$ by an element of norm one keeps the norm unchanged. It gives another pair only if both transformed coordinates belong to the original real input set.

The September adaptive construction produced the exact finite $(N,|Q(A,A)|)$ values $(2,2)$, $(4,5)$, $(8,17)$, $(16,68)$, $(32,269)$. An all-level signed-root support bound, RC15/4, would give a uniform saving under its hypotheses. It remains open. Subsequent modular-lift and cross-block studies addressed the cost of realizing whole Cartesian products. \src{E8}

In October the balanced-factor construction reached a decisive obstruction. For that unchanged sequence of carriers $A_r$ there are fixed $M\ge1$ and $r_*$ such that at all $r\ge r_*$,
\[
    |A_r|^2/(18M)\le |Q(A_r,A_r)|\le |A_r|^2.
\]
This bounds the actual image, ruling out a fixed exponent saving for that family. The proof restricts multiplicity on a positive proportion of sources, including alternatives outside the selected moves. \src{E8}

A separate mixed-signature lattice carrier $A_j$ uses allowed coefficients inside fixed bounds at the archimedean embeddings. On the recorded tower with level parameter $d_j$, one radius $T_*$ supplies at least $2^{d_j}$ norm-one order units per level. On the $M_j=|A_j|^2-1$ nonzero source pairs define $R_j(p)$ as the number of actual pairs in the same unit orbit as $p$. Each occupied orbit contributes one to
\[
  \nu_j=\sum_p\frac{1}{R_j(p)}
       =M_j\,\mathbb E_p\!\left[1/R_j(p)\right],
  \qquad |Q(A_j,A_j)|\le 1+\nu_j.
\]
An upper bound decaying exponentially in $d_j$ for the inverse moment would give a power saving. It remains open. Counting available units does not tell how often both transformed coordinates fall back inside the set. \src{E8}

\subhead{4 November\enspace /\enspace A separate proved family}
The project reports Theorem 1 for every fixed $a,b,c,e\in\Z$. The polynomial $(ax+by)^2+cx+ey$ admits a fixed power saving on arbitrarily large finite real squares. The condition $ab(ae-bc)\ne0$ singles out nonspecial examples. The theorem includes degenerate cases. The proof imports a field-supply theorem (Proposition 3.8 in its source). The audit separates that input from internal deductions MP-01--MP-07. The quadratic part of this family has rank at most one. T03 has rank two. All four targets remain open. \src{E10}
